\begin{longtable}{|p{3.4cm}|p{10.3cm}|}
\caption{Đặc tả kịch bản thu hồi chứng thực}
\label{tab:vc-revocation-scenario} \\

\hline
\textbf{Mục} & \textbf{Nội dung} \\
\hline
\endfirsthead

\hline
\textbf{Mục} & \textbf{Nội dung} \\
\hline
\endhead

\hline
\endfoot

\hline
\endlastfoot

Mã kịch bản & KB6 \\
\hline

Tác nhân chính & Bên phát hành (thu hồi, đình chỉ) hoặc quản trị viên (thu hồi, kích hoạt lại) \\
\hline

Thành phần tham gia &
Ví của người thực hiện, mô-đun \cbox{vc\_registry}, tầng chỉ mục, dịch vụ xác minh \\
\hline

Mục đích &
Cập nhật trạng thái của một chứng thực để các lần xác minh sau đó từ chối nó, dù chữ ký vẫn hợp lệ về mặt toán học \\
\hline

Điều kiện tiên quyết &
\begin{itemize}[leftmargin=*, nosep]
    \item Chứng thực tồn tại trên chuỗi
    \item Với thao tác của bên phát hành: người thực hiện đúng là bên đã cấp chứng thực đó, và định danh của họ còn hoạt động. \textbf{Không} yêu cầu thẩm quyền còn hiệu lực
    \item Với thao tác của quản trị viên: người thực hiện là quản trị viên gốc hoặc quản trị viên đang hoạt động
\end{itemize} \\
\hline

Phân biệt bốn thao tác &
\begin{itemize}[leftmargin=*, nosep]
    \item \textbf{Đình chỉ} (bên phát hành): trạng thái tạm thời, có thể đảo ngược --- nhưng chỉ quản trị viên đảo ngược được
    \item \textbf{Kích hoạt lại} (chỉ quản trị viên): đưa chứng thực từ tạm đình chỉ về đang hoạt động
    \item \textbf{Thu hồi} (bên phát hành): tuyên bố vĩnh viễn, không thể đảo ngược
    \item \textbf{Thu hồi} (quản trị viên): dùng khi cần cưỡng chế, chẳng hạn khi bên phát hành không còn hoạt động
\end{itemize} \\
\hline

Luồng thực hiện chính &
\begin{enumerate}[leftmargin=*, nosep]
    \item Người thực hiện chọn chứng thực trong bảng danh sách và chọn thao tác
    \item Giao diện đọc số hiệu chứng thực trên chuỗi từ chỉ mục --- số này không do người dùng gõ tay, để tránh gõ nhầm một chữ số và xử lý nhầm văn bằng của người khác
    \item Ví ký giao dịch gọi hàm tương ứng kèm lý do
    \item Hợp đồng kiểm điều kiện, cập nhật trạng thái, giảm bộ đếm số đang hiệu lực, rồi phát sự kiện
    \item Tầng chỉ mục xác minh giao dịch, cập nhật trạng thái trong chỉ mục và ghi nhật ký
\end{enumerate} \\
\hline

Hai trường hợp đặc biệt được hỗ trợ &
\begin{itemize}[leftmargin=*, nosep]
    \item \textbf{Đơn vị đã hết hạn ủy quyền vẫn thu hồi được} chứng thực mình đã cấp. Đây là kịch bản quan trọng nhất của nguyên tắc bất đối xứng: phát hiện gian lận thường xảy ra sau khi chu kỳ ủy quyền đã kết thúc
    \item \textbf{Chứng thực đã hết hạn vẫn thu hồi được}, vì hết hạn và thu hồi mang hai thông điệp khác nhau: hết hạn nói chứng thực không còn giá trị từ một thời điểm, thu hồi nói nó chưa bao giờ đáng tin
\end{itemize} \\
\hline

Luồng ngoại lệ &
\begin{itemize}[leftmargin=*, nosep]
    \item Người thực hiện không phải bên đã cấp $\rightarrow$ từ chối với mã lỗi không phải bên cấp
    \item Bên phát hành cố kích hoạt lại chứng thực $\rightarrow$ từ chối, thao tác này chỉ thuộc quản trị viên
    \item Thu hồi một chứng thực đã bị thu hồi $\rightarrow$ từ chối
    \item Định danh của bên thực hiện đã bị thu hồi $\rightarrow$ từ chối
\end{itemize} \\
\hline

Kết quả sau cùng &
Trạng thái mới tra cứu được trên chuỗi; bước thứ năm của quy trình xác minh từ chối chứng thực từ thời điểm này; bộ đếm số đang hiệu lực giảm \\
\hline

Tính chất được kiểm chứng &
F11, F12, N2 --- và các hành vi T13 trong mô hình đe dọa \\
\hline

\end{longtable}
